\documentclass[11pt]{article}
% Shared preamble for Theseus user guides.
% Compiled with pdflatex. See docs/guides/BUILD.md for rationale.

\usepackage[a4paper,margin=0.85in]{geometry}
\usepackage[T1]{fontenc}
\usepackage[utf8]{inputenc}

% Body: Palatino (via mathpazo). Headings: Helvetica (sans).
% Palatino is requested in the house style as a substitute for Charter
% because Charter is not part of the standard TeX Live distribution
% that the build container ships.
\usepackage{mathpazo}
\usepackage[scaled=0.92]{helvet}
\usepackage{courier}
\renewcommand{\familydefault}{\rmdefault}

\usepackage{microtype}
\usepackage{parskip}
\usepackage{enumitem}
\usepackage{booktabs}
\usepackage{longtable}
\usepackage{titlesec}
\usepackage{xcolor}
\usepackage{fancyhdr}
\usepackage{fancyvrb}
\usepackage{graphicx}
% Allow `_` in text mode without escaping. The screenshot file names and
% some environment-variable references in the guides contain underscores.
\usepackage{underscore}
\usepackage{tcolorbox}
\tcbuselibrary{breakable,skins}
\usepackage[hidelinks,colorlinks=true,linkcolor=black,urlcolor=teal,citecolor=black]{hyperref}

% Sans for chapter and section headings; serif body keeps prose readable
% on screen and in print.
\titleformat{\section}{\sffamily\Large\bfseries}{\thesection}{0.7em}{}
\titleformat{\subsection}{\sffamily\large\bfseries}{\thesubsection}{0.6em}{}
\titleformat{\subsubsection}{\sffamily\normalsize\bfseries}{\thesubsubsection}{0.5em}{}

\pagestyle{fancy}
\fancyhf{}
\fancyhead[L]{\small\sffamily Theseus User Guide}
\fancyhead[R]{\small\sffamily\thepage}
\renewcommand{\headrulewidth}{0.3pt}

\setlength{\parindent}{0pt}

% Convenience commands shared by every guide.
\newcommand{\page}[1]{\textbf{#1}}
\newcommand{\file}[1]{\texttt{#1}}
\newcommand{\route}[1]{\texttt{#1}}
\newcommand{\env}[1]{\texttt{#1}}
\newcommand{\code}[1]{\texttt{#1}}

% Callout boxes used by every guide.
\newtcolorbox{tldr}{
  colback=teal!5!white,colframe=teal!55!black,
  fonttitle=\sffamily\bfseries,title=TL;DR,
  breakable,boxrule=0.4pt,arc=2pt
}
\newtcolorbox{youwillneed}{
  colback=gray!4!white,colframe=gray!55!black,
  fonttitle=\sffamily\bfseries,title=You will need,
  breakable,boxrule=0.4pt,arc=2pt
}
\newtcolorbox{notcovered}{
  colback=orange!5!white,colframe=orange!55!black,
  fonttitle=\sffamily\bfseries,title=This guide does NOT cover,
  breakable,boxrule=0.4pt,arc=2pt
}
\newtcolorbox{preview}{
  colback=yellow!8!white,colframe=yellow!55!black,
  fonttitle=\sffamily\bfseries,title=Preview --- partially shipped,
  breakable,boxrule=0.4pt,arc=2pt
}
\newtcolorbox{nextup}{
  colback=blue!4!white,colframe=blue!55!black,
  fonttitle=\sffamily\bfseries,title=Where to read next,
  breakable,boxrule=0.4pt,arc=2pt
}

% Insert a screenshot only when the file exists, so the build does not
% break before the Playwright capture run is wired up.
\newcommand{\screenshot}[3]{%
  \IfFileExists{screenshots/#1}{%
    \begin{figure}[h]\centering
      \includegraphics[width=\linewidth]{screenshots/#1}
      \caption{#2}\label{fig:#3}
    \end{figure}%
  }{%
    \begin{center}\small\itshape%
      [screenshot \texttt{screenshots/#1} not yet captured --- #2]%
    \end{center}%
  }%
}


\title{\sffamily\bfseries The Oracle, Research Advisor, and Reading Queue\\[2pt]\large Guide 3 of 6}
\date{}
\author{}

\begin{document}
\maketitle
\thispagestyle{empty}

\begin{tldr}
The Oracle is the firm's question-answering surface. It does not run a
free chat against a foundation model; it answers \emph{from} the corpus
the workshop has built and refuses --- or warns --- when the corpus is
too thin. This guide explains the citation contract (every quoted span
must literally appear in the source), the confidence bands, the
hallucination warning, and what to do with a low-confidence answer. It
also covers two adjacent surfaces: the \page{Research Advisor} (what
to read next, and why) and the \page{Reading Queue} (your personal
stack).
\end{tldr}

\begin{youwillneed}
\begin{itemize}[leftmargin=*,itemsep=0.2em,topsep=0.2em]
  \item A signed-in Codex session.
  \item A working corpus --- at minimum a handful of uploads and a few
        accepted principles, otherwise the Oracle will mostly abstain.
\end{itemize}
\end{youwillneed}

\begin{notcovered}
\begin{itemize}[leftmargin=*,itemsep=0.2em,topsep=0.2em]
  \item How claims, principles, and conclusions get into the corpus
        --- Guide 2.
  \item Currents opinions, which use a different (event-grounded)
        prompt path --- Guide 4.
\end{itemize}
\end{notcovered}

\section{What the Oracle is}

The Oracle lives at \route{/ask} (public, the firm's open Q\&A) and at
\route{/codex-ask} (founder-facing, with access to private material).
Both accept a question in natural language and return:

\begin{itemize}[leftmargin=*,itemsep=0.2em]
  \item an answer, written in the firm's voice;
  \item inline citations to claims, conclusions, principles, or
        uploads;
  \item a confidence band: \texttt{high}, \texttt{medium},
        \texttt{low}, or \texttt{abstain};
  \item a hallucination-risk badge when the answer text contains a
        span the citation chain does not unambiguously support.
\end{itemize}

It does not return: free-form essays not grounded in the corpus,
opinions on outside events (that is Currents' job), or speculative
projections (that is the forecasts portfolio's job).

\screenshot{03_oracle_answer.png}{An Oracle answer with inline citations, confidence band, and hallucination badge.}{or-answer}

\section{How the Oracle composes an answer}

The path under the hood is short and conservative. The Oracle:

\begin{enumerate}[leftmargin=*,itemsep=0.2em]
  \item Encodes your question into a fingerprint.
  \item Retrieves the top-ranked principles and conclusions whose
        fingerprints are nearest to the question's, filtered by a
        minimum conviction threshold.
  \item Pulls the supporting claims for each retrieved principle.
  \item Hands the principles (as axioms) and claims (as evidence) to
        a language model with a strict contract: produce an answer,
        cite which principle grounds each part, note caveats.
  \item Validates that the answer doesn't contradict the retrieved
        principles by running a quick coherence check.
  \item Generates an adversarial counter-position to the leaned-on
        principles --- the strongest case against them --- and
        includes its strongest objection as a caveat.
  \item Runs a final consistency check before returning the answer.
\end{enumerate}

The result is not a one-shot LLM completion; it is an answer that has
been compared to the firm's existing positions and challenged by its
own machinery before you see it.

\section{What counts as a citation}

A citation is a structured reference back to one of:

\begin{itemize}[leftmargin=*,itemsep=0.2em]
  \item \textbf{Claim} --- an extracted statement from an upload. The
        most granular citation. Click-through lands on the source
        span.
  \item \textbf{Conclusion} --- a synthesized firm position, possibly
        published as an article or as a reviewed conclusion. Citing
        one means the Oracle is leaning on a vetted firm output, not
        a raw extracted claim.
  \item \textbf{Principle} --- a distilled, reusable position. The
        Oracle treats principles as soft constraints, not hard rules.
  \item \textbf{Upload} --- a direct quote or paraphrase from an
        uploaded document, attached with offset and page number.
\end{itemize}

Every citation is clickable. If a citation does not click through,
the Oracle has produced a broken reference --- treat it as a
hallucination warning even if no explicit badge is shown, and report
it through \page{Critiques}.

\section{The verbatim-citation rule}

The single load-bearing rule of the Oracle (and of every other
public-facing surface the workshop produces) is the \emph{verbatim
citation contract}: any quoted span the answer leans on must appear
\emph{word-for-word} in the cited source. The workshop's validator
checks this on every answer. When a quoted span doesn't match, the
answer is not retried with a fix-it prompt --- it is rejected, and
the Oracle abstains.

This is the firm's standing answer to fabrication. You will see it
echoed throughout the stack: Currents abstains rather than write an
opinion whose quote doesn't appear in any source; Forecasts abstains
rather than publish a prediction whose reasoning cites a span that
isn't there.

\section{Confidence bands}

The Oracle assigns each answer a band based on the strength and
breadth of the citation chain:

\begin{itemize}[leftmargin=*,itemsep=0.2em]
  \item \texttt{high} --- the answer is supported by multiple
        independent citations, including at least one conclusion or
        principle.
  \item \texttt{medium} --- the answer is supported by at least one
        citation that directly addresses the question.
  \item \texttt{low} --- citations exist but only weakly support the
        answer. Treat as a starting point, not as a finished
        position.
  \item \texttt{abstain} --- the corpus is too thin or off-topic to
        answer. The Oracle returns no answer text and explains which
        material it would have needed.
\end{itemize}

Abstention is the correct outcome for a sparse corpus. It is not a
failure. The Oracle was deliberately designed to refuse where it
would otherwise fabricate.

\section{The hallucination warning}

A hallucination badge appears when post-hoc verification finds the
answer text contains a span that the cited sources do not
unambiguously support. The badge does not necessarily mean the answer
is wrong; it means the support is weaker than the prose makes it
sound. Standard response:

\begin{enumerate}[leftmargin=*,itemsep=0.2em]
  \item Read the cited sources. Decide whether the unsupported span
        is a paraphrase that survives or an actual overreach.
  \item If it is an overreach, rephrase the question to be narrower
        and re-ask. The Oracle's narrower answer is usually correct.
  \item If you find the Oracle consistently overreaching on a topic,
        log a critique --- the critiques pipeline feeds back into
        prompt and policy tuning.
\end{enumerate}

\section{What to do with low-confidence answers}

Low confidence does not mean ``probably wrong.'' It means ``the corpus
will not let me commit to a sharper answer.'' Useful responses:

\begin{itemize}[leftmargin=*,itemsep=0.2em]
  \item \textbf{Add corpus.} Upload one or two documents that directly
        address the question. Re-ask. The same question with denser
        material usually moves up a band.
  \item \textbf{Narrow the question.} A low-confidence answer to a
        broad question often becomes a high-confidence answer to a
        more specific sub-question.
  \item \textbf{Promote a principle.} If the topic recurs and the
        firm has a working position, draft and accept a principle
        (Guide 2). The Oracle will start citing it.
  \item \textbf{Accept abstention.} For genuinely outside-corpus
        questions, the right answer is ``the firm does not have a
        position on this'' --- the Oracle will say so.
\end{itemize}

\section{The Research Advisor}

\page{Research} (\route{/research}) is the other side of the Oracle:
instead of answering a question, it tells you what the firm should be
reading next.

It's driven by a small set of inputs:

\begin{itemize}[leftmargin=*,itemsep=0.2em]
  \item The firm's recent abstentions. If the Oracle keeps refusing
        on a topic, that topic is a candidate for a reading
        recommendation.
  \item The firm's open contradictions. Two claims pulling against
        each other often have a resolving source somewhere.
  \item Drift events --- principles whose meaning has been shifting.
  \item The recent forecasts that have resolved against the firm
        (Brier score worse than baseline). A bad forecast often
        signals a methodology gap.
\end{itemize}

Each recommendation is a small block: a suggested title, a one-
paragraph summary of why the firm should read it, a list of reading
URIs (papers, articles, podcasts), and which upload or which open
question prompted the suggestion.

\section{The Reading Queue}

\page{Reading queue} (\route{/reading-queue}) is your personal stack
of suggestions. You can:

\begin{itemize}[leftmargin=*,itemsep=0.2em]
  \item Promote a research suggestion onto your queue.
  \item Mark items \texttt{reading}, \texttt{read}, or
        \texttt{deferred}.
  \item Attach quick notes that the workshop reads if you later upload
        the item.
\end{itemize}

The Reading Queue is intentionally low-ceremony. It's a list of things
to read with light state, not a project tracker.

\section{Inverse inference --- ``what does this event mean for our corpus?''}

A specialty mode of the Oracle, available from the question form,
answers the inverse question: given that some event happened, what
does it \emph{imply} about (or refute in) the firm's existing
positions? It returns:

\begin{itemize}[leftmargin=*,itemsep=0.2em]
  \item A list of \textbf{supporting implications} --- existing
        principles or conclusions that this event strengthens.
  \item A list of \textbf{refuted implications} --- principles or
        conclusions this event undermines, with severity (mild,
        moderate, or severe).
  \item A list of \textbf{irrelevant} principles, for transparency.
  \item A \textbf{blindspot report} --- entities, mechanisms, or
        adjacent topics the firm's corpus is missing in light of the
        event.
\end{itemize}

Severe refutations show up on the unified attention queue (Guide 6).
Blindspots show up on \page{Research} as new reading suggestions.

\section{What the Oracle will not do}

\begin{itemize}[leftmargin=*,itemsep=0.2em]
  \item It will not answer from general world knowledge alone. If the
        corpus does not contain relevant material, it abstains.
  \item It will not generate predictions, market opinions, or trade
        suggestions. Those routes are gated behind the operator
        confirmation flow described in Guide 5.
  \item It will not assert a firm position that contradicts an
        accepted principle without an explicit flag in the answer
        signalling the conflict.
  \item It will not surface private uploads in answers visible to
        readers outside the firm. The public Oracle at \route{/ask}
        is gated to the public corpus; \route{/codex-ask} (signed in)
        sees the full corpus available to you.
\end{itemize}

\section{Where the docs are}

\begin{itemize}[leftmargin=*,itemsep=0.2em]
  \item \file{docs/methods/MQS\_Specification.md} --- citation
        contract.
  \item The other five guides in this series.
\end{itemize}

\begin{nextup}
Continue with \href{04_Currents.pdf}{Guide 4 --- Currents} for the
firm's outward-facing live commentary surface, which uses a related
but distinct citation contract.
\end{nextup}

\end{document}
